\chapter{Keliling, Luas, Volume}\label{ch:g6:measure}

Berapa panjang pagarnya, seberapa besar sawahnya, berapa banyak air
yang muat di dalam tangki? Tiga pertanyaan yang berbeda, tiga besaran
yang berbeda --- \omterm{def:g6:measure:perimeter}{keliling}, \omterm{def:g6:measure:area}{luas}, \omterm{def:g6:measure:volume}{volume} --- masing-masing dengan
satuannya sendiri. Mencampuradukkannya adalah kesalahan yang paling
sering terjadi dalam geometri; bab ini memilahnya untuk selamanya.

\section{Panjang dan keliling}

\begin{definition}[Keliling]\label{def:g6:measure:perimeter}
\emph{Keliling}\index{keliling} sebuah bangun adalah panjang seluruh
tepinya. Ia diukur dengan satuan panjang: milimeter (mm), sentimeter
(cm), meter (m), kilometer (km), dengan
\[
1 \text{ km} = 1000 \text{ m}, \qquad
1 \text{ m} = 100 \text{ cm}, \qquad
1 \text{ cm} = 10 \text{ mm}.
\]
\end{definition}

\begin{example}\label{ex:g6:measure:perimeter}
Persegi panjang dengan panjang $L$ dan lebar $w$ punya \omterm{def:g6:measure:perimeter}{keliling}
\[
P = L + w + L + w = 2 \times (L + w).
\]
Untuk persegi panjang $7$ cm kali $4$ cm: $P = 2 \times (7 + 4) = 2 \times 11 =
22$ cm. Persegi dengan sisi $c$ punya \omterm{def:g6:measure:perimeter}{keliling} $4c$.
\end{example}

\begin{proposition}[Lingkar sebuah lingkaran]\label{prop:g6:measure:circle}
\omterm{def:g6:measure:perimeter}{Keliling} (atau \emph{lingkar}) sebuah lingkaran berjari-jari $r$ adalah
\[
P = 2 \pi r ,
\]
dengan $\pi \approx 3.14$ adalah bilangan yang sama untuk setiap lingkaran.
\end{proposition}
\admitted

\begin{example}
Sebuah kolam berbentuk lingkaran berjari-jari $5$ m. Tepinya berukuran
$2 \times \pi \times 5 = 10\pi \approx 31.4$ m. Simpanlah nilai
tepatnya $10\pi$ selama mungkin; \omterm{def:g4:numbers:round}{bulatkan} hanya pada akhirnya.
\end{example}

\section{Luas}

\begin{definition}[Luas]\label{def:g6:measure:area}
\emph{Luas}\index{luas} sebuah bangun mengukur permukaan yang
ditutupinya: berapa persegi satuan yang muat di dalamnya. Satuannya:
cm$^2$ (persegi dengan sisi $1$ cm), m$^2$, km$^2$, \dots\ Hati-hati:
\[
1 \text{ m}^2 = 100 \times 100 \text{ cm}^2 = 10\,000 \text{ cm}^2
\]
--- satu meter persegi adalah persegi $100$ cm kali $100$ cm, jadi
setiap anak tangga satuannya bernilai $100$, bukan $10$.
\end{definition}

\begin{omfigure}
\begin{tikzpicture}[scale=0.68]
  \draw[gray!50, thin] (0,0) grid (7,4);
  \draw[very thick, omDef] (0,0) rectangle (7,4);
  \fill[omThm!25] (0,3) rectangle (1,4);
  \draw[omThm, thick] (0,3) rectangle (1,4);
  \node[font=\small, omThm] at (2.6,3.5) {$1$ persegi satuan};
  \node[font=\small] at (3.5,-0.5) {$7 \times 4 = 28$ persegi satuan};
\end{tikzpicture}

{\small \omterm{def:g6:measure:area}{Luas} sebuah persegi panjang: $7$ kolom yang masing-masing
berisi $4$ persegi satuan, jadi seluruhnya $7 \times 4 = 28$ persegi.
Inilah sebabnya \omterm{def:g6:measure:area}{luas} $=$ panjang $\times$ lebar.}
\end{omfigure}

\begin{proposition}[Rumus dasar luas]\label{prop:g6:measure:formulas}
\begin{center}
\renewcommand{\arraystretch}{1.5}
\begin{tabular}{l|l}
bangun & \omterm{def:g6:measure:area}{luas} \\
\hline
persegi panjang ($L \times w$) & $A = L \times w$ \\
persegi (sisi $c$) & $A = c^2$ \\[4pt]
\omterm{def:g6:shapes:triangles}{segitiga siku-siku} (sisi siku-siku $a$, $b$) & $A = \dfrac{a \times b}{2}$ \\[4pt]
cakram (\omterm{def:g3:shapes:shapes}{jari-jari} $r$) & $A = \pi r^2$ \\
\end{tabular}
\end{center}
\end{proposition}

\begin{proof}[Pembuktian untuk segitiga siku-siku]
Dua salinan \omterm{def:g6:shapes:triangles}{segitiga siku-siku} dengan sisi siku-siku $a$ dan $b$, yang
direkatkan sepanjang sisi miringnya, membentuk persegi panjang
$a \times b$. Jadi \omterm{def:g6:measure:area}{luas} segitiganya adalah setengah \omterm{def:g6:measure:area}{luas} persegi
panjang itu: $\frac{ab}{2}$. (Rumus persegi panjangnya adalah
pembilangan persegi satuan di atas; rumus cakramnya diterima tanpa
bukti, lihat \cref{ch:g7:areas}.)
\end{proof}

\begin{omfigure}
\begin{tikzpicture}[scale=0.75]
  \draw[thick, fill=omDef!15] (0,0) -- (4,0) -- (0,2.5) -- cycle;
  \draw[thick, omExo, dashed] (4,0) -- (4,2.5) -- (0,2.5);
  \draw[thin] (0.3,0) -- (0.3,0.3) -- (0,0.3);
  \node[below, font=\small] at (2,-0.1) {$a$};
  \node[left, font=\small] at (0,1.25) {$b$};
  \node[font=\small, omExo] at (2.8,1.7) {salinan kedua};
\end{tikzpicture}

{\small \omterm{def:g6:shapes:triangles}{Segitiga siku-siku} adalah setengah persegi panjang: dari sana
datang rumus $\frac{a\times b}{2}$.}
\end{omfigure}

\begin{example}[Bangun gabungan]\label{ex:g6:measure:composite}
Sebuah ruangan berbentuk L adalah persegi panjang $6$ m $\times$ $4$ m
dengan pojok persegi $2$ m $\times$ $2$ m yang dibuang. Luasnya,
selangkah demi selangkah:
\begin{enumerate}
  \item seluruh persegi panjang: $6 \times 4 = 24$ m$^2$;
  \item persegi yang dibuang: $2 \times 2 = 4$ m$^2$;
  \item \omterm{def:g6:measure:area}{luas} yang tersisa: $24 - 4 = 20$ m$^2$.
\end{enumerate}
\emph{\omterm{def:g6:measure:perimeter}{Kelilingnya}} bukan $20$ apa pun: berjalan mengelilingi huruf L
itu, tepinya tetap berukuran $6 + 4 + 6 + 4 = 20$ m --- kedua potongan
pojoknya menggantikan dua bagian dinding yang sama panjang. Bilangannya
sama secara kebetulan, tetapi yang satu meter persegi, yang lain meter!
\end{example}

\begin{remark}[Keliling sama, luas berbeda]
Dua bangun dapat punya \omterm{def:g6:measure:perimeter}{keliling} yang sama dan \omterm{def:g6:measure:area}{luas} yang sangat
berbeda: persegi $5 \times 5$ dan persegi panjang $9 \times 1$
sama-sama berkeliling $20$, tetapi luasnya $25$ dan $9$. \omterm{def:g6:measure:perimeter}{Keliling} dan
\omterm{def:g6:measure:area}{luas} benar-benar besaran yang saling bebas.
\end{remark}

\section{Volume}

\begin{definition}[Volume]\label{def:g6:measure:volume}
\emph{Volume}\index{volume} sebuah \omterm{def:g5:solids:def}{bangun ruang} mengukur ruang yang
diisinya: berapa \omterm{def:g5:solids:def}{kubus} satuan yang muat di dalamnya. Satuannya: cm$^3$,
m$^3$, \dots\ dengan $1$ m$^3 = 1\,000\,000$ cm$^3$ (setiap anak tangga
bernilai $1000$). Untuk zat cair orang juga memakai \emph{liter}:
\[
1 \text{ L} = 1 \text{ dm}^3 = 1000 \text{ cm}^3,
\qquad
1 \text{ m}^3 = 1000 \text{ L}.
\]
\end{definition}

\begin{proposition}[Volume sebuah balok]\label{prop:g6:measure:box}
Sebuah \omterm{def:g5:solids:def}{balok} (\emph{prisma persegi panjang}) dengan panjang $L$, lebar
$w$ dan tinggi $h$ punya \omterm{def:g6:measure:volume}{volume}
\[
V = L \times w \times h .
\]
\end{proposition}

\begin{proof}
Lapisan dasarnya memuat $L \times w$ \omterm{def:g5:solids:def}{kubus} satuan (gambar
\cref{def:g6:measure:area}, dengan \omterm{def:g5:solids:def}{kubus}), dan ada $h$ lapisan seperti
itu.
\end{proof}

\begin{omfigure}
\begin{tikzpicture}[scale=0.6]
  \foreach \y in {0,1} {
    \foreach \x in {0,...,3} {
      \foreach \z in {0,1} {
        \begin{scope}[shift={({\x+0.45*\z},{\y+0.3*\z})}]
          \draw[fill=omDef!12] (0,0) rectangle (1,1);
          \draw[fill=omDef!20] (0,1) -- (0.45,1.3) -- (1.45,1.3) -- (1,1)
            -- cycle;
          \draw[fill=omDef!28] (1,0) -- (1.45,0.3) -- (1.45,1.3) -- (1,1)
            -- cycle;
        \end{scope}
      }
    }
  }
  \node[below, font=\small] at (2.5,-0.3)
    {$4 \times 2 \times 2 = 16$ kubus satuan};
\end{tikzpicture}

{\small Sebuah \omterm{def:g5:solids:def}{balok} yang diisi \omterm{def:g5:solids:def}{kubus} satuan: $4 \times 2$ \omterm{def:g5:solids:def}{kubus} per
lapisan, dua lapisan.}
\end{omfigure}

\begin{example}\label{ex:g6:measure:aquarium}
Sebuah akuarium berukuran $60$ cm kali $30$ cm kali $40$ cm (tinggi). \omterm{def:g6:measure:volume}{Volumenya}:
\[
V = 60 \times 30 \times 40 = 72\,000 \text{ cm}^3 = 72 \text{ dm}^3
= 72 \text{ L}.
\]
Selangkah demi selangkah untuk mengubah satuannya: $1000$ cm$^3$
menjadi satu liter, dan $72\,000 \div 1000 = 72$.
\end{example}

\section{Latihan}

\begin{exercise}[$\star$]\label{exo:g6:measure:1}
Ubahlah: $3.5$ m menjadi cm; $420$ mm menjadi cm; $0.8$ km menjadi m;
$25\,000$ m menjadi km.
\end{exercise}

\begin{exercise}[$\star$]\label{exo:g6:measure:2}
Hitunglah \omterm{def:g6:measure:perimeter}{keliling}: persegi panjang $8$ cm $\times$ $3.5$ cm; persegi
dengan sisi $6.2$ cm; segitiga dengan sisi $5$ cm, $7$ cm dan
$9$ cm.
\end{exercise}

\begin{exercise}[$\star$]\label{exo:g6:measure:3}
Sebuah lintasan lari berbentuk lingkaran berjari-jari $50$ m. Berapa
panjang satu putaran (nilai tepatnya dengan $\pi$, lalu \omterm{def:g4:numbers:round}{dibulatkan} ke
meter)? Berapa putaran yang membuat sedikitnya $2$ km?
\end{exercise}

\begin{exercise}[$\star$]\label{exo:g6:measure:4}
Hitunglah \omterm{def:g6:measure:area}{luas}: persegi panjang $9$ cm $\times$ $4$ cm; persegi dengan
sisi $7$ m; \omterm{def:g6:shapes:triangles}{segitiga siku-siku} dengan sisi siku-siku $6$ cm dan $10$
cm; cakram berjari-jari $3$ cm (nilai tepatnya, lalu \omterm{def:g4:numbers:round}{dibulatkan} ke cm$^2$).
\end{exercise}

\begin{exercise}[$\star$]\label{exo:g6:measure:5}
Ubahlah: $3$ m$^2$ menjadi cm$^2$; $45\,000$ cm$^2$ menjadi m$^2$;
$2.5$ L menjadi cm$^3$; $4\,500$ L menjadi m$^3$.
\end{exercise}

\begin{exercise}[$\star$]\label{exo:g6:measure:6}
Sebuah sawah berbentuk persegi panjang, panjangnya $120$ m dan
lebarnya $85$ m. Berapa meter pagar yang diperlukan untuk mengelilinginya? Berapa luasnya?
\end{exercise}

\begin{exercise}[$\star$]\label{exo:g6:measure:7}
Hitunglah \omterm{def:g6:measure:volume}{volume} \omterm{def:g5:solids:def}{balok} $5$ cm $\times$ $4$ cm $\times$ $10$ cm, dan
\omterm{def:g6:measure:volume}{volume} \omterm{def:g5:solids:def}{kubus} dengan \omterm{def:g5:solids:def}{rusuk} $3$ cm.
\end{exercise}

\begin{exercise}[$\star$]\label{exo:g6:measure:8}
Gambarlah dua persegi panjang berbeda yang berkeliling $16$ cm, lalu
hitunglah luasnya. Persegi panjangmu yang mana yang lebih besar luasnya?
\end{exercise}

\begin{exercise}[$\star\star$]\label{exo:g6:measure:9}
Sebuah bangun berbentuk huruf T tersusun dari persegi panjang mendatar
$10 \times 2$ di atas persegi panjang tegak $2 \times 6$ (ukurannya
dalam cm). Hitunglah luasnya, lalu \omterm{def:g6:measure:perimeter}{kelilingnya} (telusurilah tepinya
dengan teliti, sambil menjumlahkan setiap sisinya).
\end{exercise}

\begin{exercise}[$\star\star$]\label{exo:g6:measure:10}
Sebuah kolam renang berbentuk \omterm{def:g5:solids:def}{balok} sepanjang $25$ m, selebar $10$ m
dan sedalam $2$ m.
\begin{enumerate}
  \item Berapa meter kubik air yang dimuatnya ketika penuh?
  \item Berapa liter itu?
  \item Kolamnya diisi $50\,000$ L setiap jam. Berapa lama pengisiannya?
\end{enumerate}
\end{exercise}

\begin{exercise}[$\star\star$]\label{exo:g6:measure:11}
Sebuah kebun berbentuk persegi dengan sisi $20$ m memuat kolam
berbentuk lingkaran berjari-jari $3$ m. Berapa \omterm{def:g6:measure:area}{luas} rumputnya (nilai
tepatnya, lalu \omterm{def:g4:numbers:round}{dibulatkan} ke m$^2$)?
\end{exercise}

\begin{exercise}[$\star\star\star$]\label{exo:g6:measure:12}
Sebatang cokelat berukuran $15$ cm $\times$ $6$ cm $\times$ $1$ cm.
Pembuatnya melipatduakan \emph{ketiga} ukurannya.
\begin{enumerate}
  \item \omterm{def:g6:measure:volume}{Volumenya} dikalikan berapa? Periksalah dengan menghitung kedua
        \omterm{def:g6:measure:volume}{volumenya}.
  \item Harganya dikalikan $4$. Apakah batang yang besar lebih
        menguntungkan daripada yang kecil?
\end{enumerate}
\end{exercise}

\section{Soal: Pagar Ratu Dido}

\begin{problem}[{Soal akhir pekan --- dengan panjang pagar yang tetap,
bangun mana yang melingkupi tanah paling luas? Persegi mengalahkan
setiap persegi panjang, lingkaran mengalahkan persegi, dan sebuah
dinding lumbung mengubah segalanya}]
\label{pb:g6:measure:1}
Legenda mengisahkan bahwa Ratu Dido, ketika mendarat di pesisir
Afrika, diberi ``tanah seluas yang dapat dilingkupi sehelai kulit
sapi'' --- maka ia memotong kulit itu menjadi satu jalur yang sangat
panjang dan tipis lalu melingkupi tanah yang cukup untuk mendirikan
kota Kartago. Soalnya kini menjadi soalmu: seorang petani punya tepat
$20$ m pagar. \Cref{exo:g6:measure:8} memperlihatkan bahwa dua kandang
dengan \omterm{def:g6:measure:perimeter}{keliling} yang sama dapat punya \omterm{def:g6:measure:area}{luas} yang berbeda; soal ini
mencari kandang yang \emph{terbaik} --- dan menemukan bahwa jawabannya
berubah sama sekali ketika dinding lumbung meminjamkan satu sisi
secara cuma-cuma.

\medskip
\noindent\textbf{Bagian I --- Dua puluh meter pagar.}
\begin{enumerate}
  \item Kandangnya harus berbentuk persegi panjang yang memakai
        seluruh $20$ m pagar itu. Periksalah bahwa kandang
        $1 \times 9$ dan kandang $2 \times 8$ sama-sama memenuhi
        syarat, lalu hitunglah luasnya.
  \item Jelaskan mengapa panjang dan lebar setiap kandang yang
        memenuhi syarat berjumlah $10$ m. Lalu buatlah tabel lengkap
        kandang dengan sisi bilangan bulat (dari $1 \times 9$ sampai
        $5 \times 5$) beserta luasnya. Mana yang terbaik?
  \item Apakah sisi desimal layak dicoba? Hitunglah \omterm{def:g6:measure:area}{luas} kandang
        $4.5 \times 5.5$ dan kandang $4.9 \times 5.1$, lalu
        bandingkan dengan persegi $5 \times 5$. Apa dugaanmu?
  \item Pertanyaan 2 mengubah soal pagar itu menjadi soal bilangan
        yang murni: \emph{di antara semua pasangan bilangan yang
        berjumlah $10$, pasangan mana yang \omterm{def:g2:mult:def}{hasil kalinya} terbesar?}
        Jawablah dari tabelmu, lalu nyatakan aturan umum yang kamu
        amati.
  \item Inilah alasan mengapa menjauh dari bentuk persegi selalu
        merugi, dengan gunting alih-alih aljabar: mulailah dari
        kandang $5 \times 5$ lalu ubahlah menjadi kandang
        $6 \times 4$ dengan membuang satu jalur dan merekatkan jalur
        lain sebagai gantinya. Jalur mana yang dibuang, jalur mana
        yang ditambahkan, dan mengapa penukaran itu kehilangan tepat
        satu meter persegi? Jelaskan mengapa langkah berikutnya (ke
        $7 \times 3$) kehilangan lebih banyak lagi.
\end{enumerate}

\medskip
\noindent\textbf{Bagian II --- Lumbungnya, dan lingkarannya.}
\begin{enumerate}[resume]
  \item Kandangnya sekarang dibangun menempel pada dinding lumbung
        yang panjang: dindingnya menggantikan satu sisi panjang
        kandang, dan $20$ m pagarnya hanya menutupi tiga sisi yang
        lain. Buatlah tabel kandang dengan ukuran bilangan bulat
        (lebar $w$, panjang $L$ sepanjang dinding, $w + L + w = 20$)
        beserta luasnya. Kandang mana yang menang sekarang --- dan
        apakah ia persegi?
  \item Jelaskan pemenangnya dengan sebuah cermin
        (\cref{ch:g6:symmetry}): cerminkan kandangnya terhadap
        dinding lumbung lalu perhatikan kandang gandanya. Berapa
        pagar yang dipakai kandang ganda itu, kandang ganda mana yang
        terbaik menurut Bagian I, dan itu menjadikan kandang aslinya
        apa?
  \item Dido tidak membangun persegi panjang. Lengkungkan $20$ m
        pagar itu menjadi lingkaran: dengan $\pi \approx 3.14$,
        hitunglah \omterm{def:g3:shapes:shapes}{jari-jarinya} (sampai cm), lalu luasnya (sampai
        m$^2$) (\cref{prop:g6:measure:circle},
        \cref{prop:g6:measure:formulas}). Bandingkan dengan persegi
        $5 \times 5$.
  \item Soal kebalikannya: diinginkan kandang seluas tepat
        $36$ m$^2$, dengan pagar sesedikit mungkin. Bandingkan
        persegi panjang $1 \times 36$, $2 \times 18$, $3 \times 12$,
        $4 \times 9$ dan $6 \times 6$: berapa \omterm{def:g6:measure:perimeter}{kelilingnya}? Bangun
        mana yang paling murah, dan bagaimana pertanyaan ini menjadi
        bayangan cermin Bagian I?
  \item Dalam satu kalimat: mengapa kaleng, pipa dan tangki air
        begitu sering \emph{bulat}?
\end{enumerate}

\medskip
\noindent\textbf{Bagian III --- \omterm{def:g6:measure:perimeter}{Keliling} dan \omterm{def:g6:measure:area}{luas} itu orang asing.}
\begin{enumerate}[resume]
  \item Carilah persegi panjang yang \omterm{def:g6:measure:perimeter}{kelilingnya} melebihi $100$ m
        tetapi luasnya kurang dari $1$ m$^2$. (Sangat panjang dan
        sangat tipis. \omterm{def:g6:decimals:places}{Bilangan desimal} diperbolehkan.)
  \item Benar atau salah: ``dari dua bangun, yang \omterm{def:g6:measure:perimeter}{kelilingnya} lebih
        besar luasnya juga lebih besar.'' Berikan contoh penyangkal
        dari soal ini.
  \item Ambillah persegi panjang $6 \times 4$ lalu potonglah takik
        persegi $1 \times 1$ di tengah salah satu sisi panjangnya
        (takiknya menganga ke luar, seperti gigi yang hilang).
        Hitunglah \omterm{def:g6:measure:area}{luas} dan \omterm{def:g6:measure:perimeter}{keliling} bangun bertakik itu, lalu
        bandingkan keduanya dengan bangun aslinya. Apa yang terjadi
        pada masing-masing?
  \item Pembuat peta mengetahui fakta yang aneh: makin rinci petanya,
        makin \emph{panjang} garis pantainya terukur --- setiap
        perbesaran memunculkan takik kecil yang baru, dan pertanyaan
        13 memperlihatkan apa yang dikerjakan setiap takik. Jelaskan
        dalam satu atau dua kalimat mengapa ``panjang garis pantai
        Bretagne'' adalah bilangan yang licin, sedangkan ``\omterm{def:g6:measure:area}{luas}
        Bretagne'' tidak.
  \item Ujian sang petani. Dengan $36$ m pagar dan dinding lumbung
        yang tersedia, carilah kandang persegi panjang yang terbaik
        (pakailah cermin pada pertanyaan 7), sebutkan luasnya, lalu
        bandingkan dengan kandang terbaik yang dibangun dengan $36$ m
        di lapangan terbuka. Berapa banyak yang diperoleh petani itu
        berkat dinding lumbungnya?

\end{enumerate}
\end{problem}
